\chapter{État de l'art}

\section{Journalisation système et réseau}

La journalisation désigne l'enregistrement des événements générés par un système, une application ou un équipement réseau. Ces journaux peuvent contenir des informations sur les authentifications, les erreurs, les accès aux fichiers, les requêtes web, les modifications de configuration ou encore les alertes de sécurité. Le guide NIST SP 800-92 rappelle d'ailleurs que la collecte, la conservation, l'analyse et la protection de ces journaux constituent des éléments importants d'une stratégie de sécurité \cite{nist80092}.

Dans la pratique, ces données se présentent sous des formats très différents. Le protocole syslog organise les messages système autour d'informations comme l'horodatage, l'hôte, la sévérité et le contenu du message \cite{rfc5424}. Windows Event Log utilise pour sa part un format propre aux environnements Microsoft. Les serveurs Apache produisent principalement des journaux textuels liés aux accès et aux erreurs, tandis que des plateformes HIDS/SIEM telles que Wazuh enrichissent les événements avec des règles, des niveaux de sévérité, des identifiants d'agents ou encore des groupes. À l'inverse, des jeux de données réseau comme CICIDS2017, UNSW-NB15 ou CIC-DDoS2019 représentent généralement les communications sous forme de flux tabulaires caractérisés par plusieurs variables numériques.

Cette diversité rend difficile l'application directe d'une méthode unique de détection. Avant de chercher à identifier un comportement anormal, il est donc nécessaire de comprendre la structure des données, d'en extraire les informations utiles et de les rapprocher dans une représentation commune.

\section{Détection d'intrusions et détection d'anomalies}

La détection d'intrusions cherche à identifier des activités malveillantes, suspectes ou non autorisées au sein d'un système informatique. Denning a proposé l'un des premiers modèles structurés reposant sur l'analyse des traces d'audit et sur l'identification des écarts par rapport à un comportement considéré comme normal \cite{denning1987intrusion}. Axelsson a ensuite proposé une taxonomie des systèmes de détection d'intrusions, en distinguant notamment les méthodes fondées sur les signatures de celles reposant sur la détection d'anomalies \cite{axelsson2000intrusion}.

La détection d'anomalies adopte une logique légèrement différente. Elle cherche avant tout à identifier des observations rares, inhabituelles ou suffisamment éloignées d'un comportement de référence. Chandola et al. distinguent notamment les anomalies ponctuelles, contextuelles et collectives \cite{chandola2009anomaly}. Dans un journal informatique, une anomalie peut ainsi prendre la forme d'une erreur isolée, d'une succession inhabituelle d'événements, d'une activité rare associée à un utilisateur ou encore d'une augmentation anormale du nombre d'événements sur une courte période.

Cette distinction est importante dans le cadre de ce mémoire. Un événement identifié comme anormal par un modèle ne constitue pas automatiquement une intrusion. Il s'agit d'abord d'un signal qui doit être replacé dans son contexte avant de pouvoir être interprété comme un incident de sécurité.

\section{Approches classiques}

Les méthodes classiques de détection reposent principalement sur des règles, des signatures, des seuils ou des statistiques relativement simples. Elles restent largement utilisées, notamment parce que leur fonctionnement est généralement compréhensible et que leurs décisions peuvent être expliquées plus facilement.

Un outil comme fail2ban peut, par exemple, surveiller des échecs répétés d'authentification et déclencher le bannissement temporaire d'une adresse IP lorsque certaines conditions sont réunies. Des plateformes HIDS/SIEM telles que OSSEC ou Wazuh utilisent quant à elles des décodeurs, des règles et différents mécanismes de corrélation afin d'identifier des situations considérées comme suspectes et de générer des alertes.

Ces approches ont une réelle valeur opérationnelle, notamment lorsque les comportements recherchés sont connus. Elles nécessitent cependant une configuration adaptée au contexte et peuvent générer des faux positifs. Leur efficacité diminue également lorsqu'un comportement nouveau n'est couvert par aucune règle ou signature existante.

Le positionnement de \logminer{} n'est donc pas celui d'un remplacement de ces outils. Le prototype est plutôt envisagé comme un complément analytique capable d'ajouter des mécanismes de détection fondés sur les données tout en conservant une validation humaine.

\section{Apprentissage automatique pour la détection}

L'apprentissage automatique permet de construire des modèles capables d'apprendre des relations ou des frontières de décision à partir des données disponibles. Dans un contexte de cybersécurité, ces méthodes peuvent être utilisées aussi bien pour la classification d'événements que pour l'identification de comportements atypiques.

Les approches supervisées, comme Random Forest, supposent la présence de données labellisées. Elles peuvent obtenir de très bonnes performances sur des données tabulaires, mais leur efficacité dépend fortement de la qualité des labels, de la manière dont les données sont séparées entre entraînement et test ainsi que de la stabilité de leur distribution.

Les méthodes non supervisées, comme Isolation Forest, fonctionnent sans labels explicites. Elles cherchent plutôt à identifier des observations rares ou éloignées du comportement majoritaire. Dans ce cas, les sorties du modèle doivent être interprétées comme des anomalies candidates et non comme des attaques confirmées.

Buczak et Guven présentent un panorama de nombreuses méthodes d'apprentissage automatique appliquées à la cybersécurité \cite{buczak2016survey}. Khraisat et al. insistent notamment sur l'importance des jeux de données utilisés, des protocoles d'évaluation et des différences entre les environnements expérimentaux et les systèmes réels \cite{khraisat2019survey}. Sommer et Paxson soulignent également les risques liés au passage direct de modèles entraînés dans des environnements expérimentaux vers des réseaux réels \cite{sommer2010outside}. Ces remarques conduisent à interpréter avec prudence les performances obtenues dans le cadre de ce mémoire.

Random Forest occupe une place importante dans les expérimentations supervisées réalisées ici. Ce modèle combine plusieurs arbres de décision afin de limiter la variance et d'améliorer la robustesse de la prédiction \cite{breiman2001random}. Il convient particulièrement bien aux données tabulaires, comme celles issues de flux réseau ou de journaux déjà transformés en caractéristiques numériques.

Son intérêt réside aussi dans sa relative simplicité d'utilisation. Il nécessite généralement moins de transformations que certains modèles linéaires et peut fournir de bonnes performances avec un effort de paramétrage raisonnable. Cela ne garantit cependant pas sa capacité à généraliser. Un modèle peut apprendre des particularités propres au jeu de données utilisé, voire exploiter indirectement certains biais introduits lors de la collecte ou de la préparation des données.

Isolation Forest est utilisé pour une partie des traitements non supervisés. Son principe repose sur l'idée que les observations inhabituelles peuvent souvent être isolées avec moins de séparations que les observations appartenant aux comportements dominants \cite{liu2008isolation}. D'autres techniques peuvent également être utilisées pour ce type d'analyse, notamment Local Outlier Factor, One-Class SVM ou encore k-means comme méthode de regroupement de référence \cite{breunig2000lof,scholkopf2001estimating,macqueen1967kmeans}.

Ces méthodes ne détectent cependant pas toutes le même type d'anomalie. Certaines analysent principalement la densité locale, d'autres construisent une frontière globale ou s'appuient sur la distance par rapport à un centre. Le choix d'une méthode ne peut donc pas reposer uniquement sur une justification théorique ; il doit également être confronté aux données et au protocole expérimental.

\section{Apprentissage profond et analyse des journaux}

L'apprentissage profond a progressivement pris une place importante dans l'analyse automatique des journaux, notamment grâce aux modèles capables d'exploiter les séquences d'événements.

DeepLog utilise par exemple des réseaux récurrents afin d'apprendre les séquences considérées comme normales et de détecter les événements qui s'écartent de ces séquences \cite{du2017deeplog}. D'autres travaux ont ensuite utilisé des modèles de messages, des représentations vectorielles et des architectures de type Transformer pour améliorer la prise en compte du contexte.

Ces approches sont particulièrement intéressantes lorsque l'ordre des événements contient une information importante. Elles nécessitent toutefois souvent davantage de données, une préparation plus précise des séquences ainsi que des ressources de calcul plus importantes.

Dans ce mémoire, \logminer{} adopte volontairement une approche plus légère. L'accent est mis sur la normalisation des journaux, le routage selon la famille de données, l'utilisation de modèles spécialisés et la restitution des résultats dans un tableau de bord auditable. L'apprentissage profond est donc étudié comme élément de comparaison et comme perspective d'évolution, mais il ne constitue pas le cœur de l'implémentation actuelle.

Les travaux sur l'analyse séquentielle des journaux montrent également que la représentation des événements joue un rôle important. LogAnomaly combine par exemple des informations séquentielles et quantitatives afin d'identifier des comportements atypiques dans des journaux non structurés \cite{meng2019loganomaly}.

De telles méthodes deviennent particulièrement pertinentes lorsque la signification d'un événement dépend fortement de ceux qui le précèdent. Leur intégration complète dans \logminer{} demanderait toutefois une évolution du pipeline actuel, notamment par l'utilisation plus systématique de fenêtres temporelles, de modèles de messages et de traitements explicitement séquentiels.

\section{Travaux récents sur les représentations des journaux}

Les recherches récentes ont progressivement rapproché l'analyse des journaux des modèles de langage et des représentations sémantiques.

LogBERT applique par exemple une approche inspirée de BERT à des séquences de journaux. Le modèle utilise un apprentissage auto-supervisé et la prédiction masquée afin de représenter le comportement normal et d'identifier les séquences inhabituelles \cite{guo2021logbert}.

SSDLog s'inscrit dans une logique semi-supervisée fondée sur une architecture enseignant--étudiant, avec pour objectif de réduire la dépendance à des jeux de données entièrement labellisés \cite{lu2023ssdlog}. LTAnomaly combine quant à lui plusieurs niveaux de représentation avec des informations sémantiques, des valeurs de composants, un LSTM et un Transformer pour analyser notamment HDFS et BGL \cite{han2023ltanomaly}.

Ces travaux montrent que les approches contemporaines ne se limitent plus aux premiers modèles récurrents. Elles cherchent à représenter conjointement la séquence, le contexte et la signification des messages.

Cette évolution ne rend pas nécessairement les modèles plus légers inutiles. Elle permet plutôt de mieux définir leur domaine d'application. \logminer{} ne cherche pas à concurrencer directement les architectures profondes spécialisées sur les benchmarks de journaux système. Son apport se situe principalement dans l'intégration de plusieurs étapes dans une chaîne cohérente : parsing, routage, détection, corrélation, audit et visualisation.

Dans cette perspective, les expériences menées sur HDFS et BGL avec Drain3 et des fenêtres temporelles doivent être considérées comme une amélioration du protocole de représentation des données, et non comme une alternative directe aux architectures Transformer spécialisées.

Les travaux récents montrent également un intérêt croissant pour le parsing automatique. LogSL traite par exemple le parsing comme une tâche auto-supervisée dans le but d'améliorer la robustesse face à la diversité des formats et au manque de labels \cite{cao2024logsl}. SemantiLog s'appuie pour sa part sur la similarité sémantique afin de rapprocher des messages même lorsque leur forme textuelle diffère \cite{shavit2024semantilog}.

Ces évolutions renforcent l'intérêt de l'utilisation de Drain3 dans \logminer{}. L'extraction de modèles de messages ne doit pas être considérée uniquement comme une opération de nettoyage préalable. La qualité du parsing influence directement les représentations qui seront ensuite utilisées par les modèles de détection.

\section{Parsing et normalisation}

Le parsing consiste à transformer les lignes brutes des journaux en informations structurées pouvant être exploitées par les traitements suivants. Des méthodes comme Drain cherchent notamment à identifier les parties stables des messages et à les regrouper en modèles \cite{he2016drain}.

Dans \logminer{}, le parsing doit prendre en charge plusieurs formats. L'objectif est d'extraire, lorsque cela est possible, un ensemble d'informations communes telles que l'horodatage, la source, l'hôte, l'utilisateur, la catégorie, le message, la sévérité ainsi que certains champs numériques.

La normalisation intervient ensuite pour rapprocher les différentes sources dans une représentation commune. Cette étape est essentielle, car sans schéma partagé, chaque format nécessiterait son propre mécanisme de traitement, de filtrage et d'affichage.

Les recherches consacrées au parsing montrent par ailleurs qu'un journal ne peut pas être réduit à une simple chaîne de caractères. Un message contient souvent une partie relativement stable, qui correspond au modèle général de l'événement, et plusieurs paramètres variables. Les travaux comparatifs sur les parseurs montrent que la manière dont cette séparation est réalisée influence les analyses effectuées ensuite \cite{zhu2019logparser}.

L'objectif retenu dans ce mémoire n'est donc pas de construire un parseur universel capable de remplacer toutes les méthodes spécialisées. L'approche consiste plutôt à combiner plusieurs parseurs contrôlés tout en conservant les informations qui ne peuvent pas être immédiatement interprétées.

\section{Jeux de données et évaluation en détection d'intrusions}

Les jeux de données occupent une place essentielle dans l'évaluation des systèmes de détection. Ils offrent un cadre commun permettant de comparer plusieurs approches et fournissent souvent des labels qui ne sont pas disponibles dans des journaux opérationnels.

CICIDS2017 a été conçu afin de proposer un trafic plus représentatif que plusieurs jeux de données historiques. Il regroupe différents scénarios d'attaque et décrit les communications à partir de caractéristiques de flux réseau \cite{sharafaldin2018cicids}. UNSW-NB15 a également été développé pour permettre l'évaluation des systèmes de détection sur un mélange de trafic normal et de plusieurs catégories d'attaques \cite{moustafa2015unsw}.

Ces jeux de données restent cependant des représentations contrôlées d'une réalité beaucoup plus large. Chaque jeu dépend de son environnement de collecte, des attaques sélectionnées, des outils employés pour générer le trafic et de la méthode utilisée pour produire les labels.

Un score obtenu sur ces données doit donc être interprété dans le cadre du protocole expérimental correspondant. Cette précaution devient particulièrement importante lorsque les performances obtenues sont proches de 1. Un score très élevé peut effectivement correspondre à un modèle performant, mais il peut également traduire une forte séparabilité des classes ou certaines particularités du jeu de données.

Les journaux système comme HDFS et BGL posent un problème supplémentaire. Dans ces données, l'anomalie peut dépendre d'une suite d'événements, d'une fenêtre temporelle ou d'une modification collective du comportement. Une analyse réalisée indépendamment ligne par ligne peut donc manquer une partie importante de l'information.

Les approches reposant sur les modèles de messages et sur les séquences peuvent être plus adaptées dans ce cas, mais elles augmentent également la complexité du pipeline. Dans ce mémoire, HDFS et BGL permettent donc à la fois d'évaluer la capacité du prototype à manipuler des journaux publics et de mettre en évidence certaines limites des représentations légères lorsque la structure temporelle ou séquentielle devient importante.

\section{Systèmes multi-agents pour la sécurité}

Les architectures multi-agents ne sont pas nouvelles dans le domaine de la détection distribuée d'intrusions. AAFID propose par exemple une organisation reposant sur plusieurs agents autonomes chargés de collecter et d'analyser des événements locaux avant leur agrégation \cite{balasubramaniyan1998aafid}.

Dans le cadre de \logminer{}, le terme \textit{agent} désigne principalement un composant logiciel spécialisé capable d'interagir avec les autres éléments de l'architecture. Ces agents ne sont pas présentés comme des systèmes cognitifs autonomes. Ils correspondent plutôt à une décomposition fonctionnelle destinée à améliorer la modularité, l'auditabilité et la possibilité de faire évoluer les différentes étapes du traitement indépendamment.

Les travaux récents consacrés à la cyberdéfense autonome insistent d'ailleurs sur le fait qu'un agent opérationnel ne peut pas être qualifié d'autonome uniquement parce qu'il exécute automatiquement certaines tâches. Des exigences supplémentaires apparaissent, notamment en matière d'explicabilité, de coordination, de gouvernance, de validation dans des environnements réalistes et de contrôle des actions \cite{vyas2026acnd}.

Les agents de \logminer{} doivent donc être replacés dans ce cadre. Ils disposent notamment d'une politique de sélection des tâches, d'une mémoire locale, d'un mécanisme de heartbeat, de plusieurs handlers spécialisés ainsi que d'une capacité de reprise de tâches via Redis Streams. En revanche, ils ne déclenchent pas automatiquement des contre-mesures sur l'infrastructure surveillée.

\section{Architectures distribuées, fédérées et robustesse temporelle}

Les recherches sur les systèmes de détection fédérés montrent un intérêt croissant pour l'apprentissage décentralisé. Cette orientation répond notamment à des contraintes de confidentialité, à la diversité des organisations participantes et au coût potentiel du transfert de grandes quantités de données brutes \cite{khraisat2024federatedids,makris2025fids}.

Ces architectures ouvrent cependant de nouvelles difficultés. Les données détenues par les différents participants ne suivent pas nécessairement la même distribution. Les échanges entre nœuds ont un coût, et le mécanisme d'agrégation peut lui-même devenir une cible d'attaque. Les questions de gouvernance, d'empoisonnement des modèles et de contrôle des participants deviennent alors importantes.

Un autre problème concerne l'évolution du comportement des systèmes au fil du temps. Les habitudes des utilisateurs changent, les applications évoluent, les volumes de données varient et de nouvelles formes d'attaque apparaissent. Cette dérive conceptuelle peut progressivement réduire la pertinence d'un modèle entraîné sur des données anciennes \cite{shyaa2024conceptdrift}.

Cette question est directement liée aux mécanismes d'apprentissage continu, au recalibrage des seuils et à l'intégration des retours humains. Dans \logminer{}, la dérive conceptuelle n'est pas encore traitée comme un module autonome. Elle est plutôt considérée comme une limite du prototype actuel et comme une perspective nécessitant un suivi temporel plus structuré, des procédures de recalibrage et une exploitation contrôlée des décisions de l'analyste.

\section{Synthèse comparative}

L'état de l'art montre qu'aucune famille de méthodes ne répond à elle seule à l'ensemble des besoins liés à l'analyse des journaux.

Les systèmes fondés sur des règles restent largement utilisés parce qu'ils sont rapides, compréhensibles et adaptés à des comportements déjà connus. Les modèles supervisés peuvent obtenir d'excellentes performances lorsque des données correctement labellisées sont disponibles. Les modèles non supervisés offrent une solution lorsque les labels sont absents, mais leurs sorties nécessitent généralement une interprétation supplémentaire. Les architectures profondes, quant à elles, permettent de mieux exploiter les séquences et le contexte, au prix d'une complexité et d'un coût de calcul plus importants.

\logminer{} se place entre ces différentes approches. Le prototype cherche surtout à organiser une chaîne complète comprenant le parsing, la normalisation, le routage vers différents modèles, la détection, la corrélation et la visualisation. Sa contribution porte donc davantage sur l'intégration et l'exploitabilité de ces différentes briques que sur la proposition d'un nouvel algorithme isolé.

\begin{table}[H]
\centering
\caption{Synthèse comparative des familles d'approches}
\label{tab:etat-art-comparaison}
\scriptsize

\begin{tabularx}{\textwidth}{p{2.8cm}p{3.5cm}p{3.5cm}X}
\toprule
\textbf{Approche} &
\textbf{Forces} &
\textbf{Limites} &
\textbf{Positionnement de \logminer{}} \\
\midrule

Règles et signatures &
Méthodes rapides, interprétables et bien adaptées aux comportements ou attaques déjà connus. &
Nécessitent une maintenance des règles et restent limitées face à certains comportements inconnus. &
Utilisées comme référence opérationnelle ; \logminer{} est présenté comme un complément et non comme un remplacement de Wazuh ou OSSEC. \\

Seuils statistiques &
Approches simples à comprendre, peu coûteuses et relativement faciles à déployer. &
Leur efficacité dépend fortement du contexte et peut se dégrader lorsque la charge ou le comportement normal évolue. &
Utilisés comme points de comparaison et complétés par des mécanismes de routage et de corrélation. \\

Apprentissage supervisé &
Peut fournir de bonnes performances lorsque des données labellisées et représentatives sont disponibles. &
Dépend de la qualité des labels et peut apprendre des particularités propres au jeu de données utilisé. &
Employé notamment sur Linux/auth, CICIDS2017 et CIC-DDoS2019. \\

Apprentissage non supervisé &
Permet d'analyser des données non labellisées et de signaler des comportements rares. &
Les résultats correspondent à des anomalies candidates et ne constituent pas directement des preuves d'intrusion. &
Utilisé avec une interprétation contextuelle et une validation par l'analyste. \\

Apprentissage profond &
Particulièrement adapté aux séquences complexes et aux grandes quantités de données. &
Demande généralement davantage de données, de calcul et d'efforts d'interprétation. &
Étudié dans l'état de l'art et retenu comme perspective plutôt que comme composant central du prototype. \\

Transformers et représentations sémantiques &
Permettent de mieux intégrer le contexte, la séquence et la proximité sémantique entre les messages. &
Demandent des protocoles séquentiels plus élaborés, davantage de ressources et une validation rigoureuse. &
Utilisés pour situer les expérimentations HDFS/BGL et orienter les évolutions futures du prototype. \\

IDS fédérés et méthodes sensibles à la dérive &
Répondent aux besoins de décentralisation, de confidentialité et d'adaptation à l'évolution temporelle. &
Soulèvent des difficultés liées aux données non-IID, à la gouvernance, à la sécurité de l'agrégation et au suivi à long terme. &
Considérés comme des prolongements possibles vers un fonctionnement multi-machine et un apprentissage continu, sans être présentés comme des capacités déjà validées. \\

\bottomrule
\end{tabularx}
\end{table}

\section{Limites identifiées dans l'état de l'art}

Plusieurs limites ressortent de cette revue et permettent de préciser le besoin auquel répond ce mémoire.

La première concerne l'hétérogénéité des données. Les journaux issus de systèmes différents ne partagent ni la même structure, ni la même granularité, ni nécessairement la même signification. Une analyse transversale devient donc difficile sans mécanisme de normalisation.

La deuxième limite concerne la disponibilité des labels. Les jeux de données publics permettent d'entraîner et d'évaluer des modèles dans un environnement contrôlé. Dans une infrastructure réelle, les événements ne sont cependant pas systématiquement labellisés, et les labels disponibles ne sont pas toujours suffisamment complets ou fiables.

La troisième limite concerne l'exploitabilité des résultats. Un modèle peut obtenir un score élevé sans nécessairement fournir les informations nécessaires à l'analyste pour comprendre pourquoi un événement a été signalé et décider de l'action à entreprendre.

Ces trois limites motivent le choix d'une approche intégratrice. Le prototype ne se concentre donc pas uniquement sur le modèle de détection. Il couvre également l'acquisition des données, leur parsing, leur normalisation, leur routage, la détection, la corrélation, la visualisation et l'audit.

Le choix est volontairement pragmatique : construire d'abord une chaîne fonctionnelle, mesurable et extensible, puis envisager l'ajout progressif de mécanismes plus complexes lorsque leur utilité peut être clairement évaluée.

\section{Typologie détaillée des journaux exploités}

Les journaux étudiés dans ce mémoire appartiennent à plusieurs familles qui ne doivent pas être considérées de la même manière.

Une première catégorie regroupe les journaux directement produits par les systèmes d'exploitation. Sous Windows, les journaux enregistrent notamment des événements associés à la sécurité, aux applications et au fonctionnement général du système. Sous Linux, les journaux d'authentification permettent par exemple d'observer les connexions, les échecs de mot de passe, certains changements d'utilisateur et différentes activités liées aux services. Ces données sont particulièrement utiles pour suivre les opérations d'administration et identifier des tentatives d'accès inhabituelles.

Une deuxième catégorie concerne les journaux applicatifs. Les journaux d'accès Apache, les fichiers d'erreurs et les traces de services web décrivent les interactions entre les utilisateurs, les applications et les ressources exposées. Ils peuvent révéler des comportements inhabituels comme une fréquence anormale de requêtes, des erreurs répétitives, des codes HTTP rares ou des tentatives d'accès à des chemins sensibles. Leur interprétation dépend toutefois fortement du contexte de l'application concernée.

Une troisième famille regroupe les événements produits par les plateformes de sécurité. Wazuh, OSSEC et d'autres solutions HIDS/SIEM ne se contentent pas de stocker des messages bruts. Elles ajoutent généralement des règles, des niveaux de sévérité, des identifiants d'agents et différentes catégories de sécurité. Ces informations sont intéressantes parce qu'elles représentent déjà une première interprétation de l'événement et permettent de comparer certains résultats du prototype avec ceux d'un outil opérationnel existant.

La quatrième catégorie concerne les flux réseau et les jeux de données tabulaires. CICIDS2017, UNSW-NB15 et CIC-DDoS2019 ne correspondent pas à des journaux textuels classiques. Les communications y sont représentées par un ensemble de caractéristiques numériques décrivant les flux.

Ces données se prêtent bien à l'utilisation de modèles supervisés. Elles ne doivent cependant pas être interprétées de la même manière que les journaux système réels. Elles fournissent surtout un cadre contrôlé permettant de mesurer les performances de certains modèles.

\section{Enjeux de normalisation}

La normalisation constitue l'un des points importants de l'architecture proposée. Elle permet de rapprocher des sources qui, au départ, ne sont pas directement comparables.

Sans représentation commune, chaque type de journal nécessiterait ses propres règles de traitement, ses propres filtres et parfois une interface spécifique. L'utilisation d'un schéma partagé permet au contraire de manipuler de manière plus homogène des informations telles que le temps, la source, la sévérité, la catégorie, l'utilisateur, le message, le score ou encore la décision associée à un événement.

Cette harmonisation ne doit toutefois pas conduire à supprimer les particularités des données d'origine. Une normalisation trop agressive peut faire disparaître des informations importantes pour l'audit ou l'interprétation d'une anomalie.

Le choix retenu dans ce travail consiste donc à utiliser un noyau de champs communs tout en conservant, dans les métadonnées, les éléments spécifiques à chaque format. Cette stratégie permet de maintenir une représentation homogène pour les traitements principaux tout en préservant les détails propres aux sources.

Elle facilite également l'évolution progressive du système. De nouveaux parseurs ou de nouveaux champs peuvent être ajoutés sans remettre en cause l'organisation générale du tableau de bord ou du pipeline de traitement.

\section{Critère d'exploitabilité pour l'analyste}

L'évaluation d'un système de détection ne peut pas reposer uniquement sur ses performances statistiques. Un outil destiné à être utilisé dans un contexte de cybersécurité doit également fournir suffisamment d'informations pour qu'un analyste puisse comprendre ce qui a été détecté.

L'utilisateur doit notamment pouvoir identifier la source de l'événement, comprendre la raison du signalement, connaître la sévérité estimée, replacer l'alerte dans son contexte temporel et consulter les décisions ou traitements déjà appliqués.

L'exploitabilité devient donc un critère de conception au même titre que la précision d'un modèle.

Cette dimension est parfois moins développée dans les travaux expérimentaux, qui se concentrent principalement sur des indicateurs comme la précision, le rappel ou le F1-score. Ces métriques restent indispensables pour mesurer les performances, mais elles ne montrent pas toujours comment le résultat sera utilisé dans une situation réelle.

Dans ce mémoire, le tableau de bord et les mécanismes d'audit sont donc intégrés directement dans l'architecture. Ils ne sont pas considérés comme de simples éléments graphiques ajoutés après l'évaluation, mais comme des composants permettant de relier la détection algorithmique à l'interprétation et à la décision humaine.
